% ============================================================
%  Поглавље 6: QAOA — Квантни алгоритам за оптимизацију
% ============================================================
\section{QAOA и MaxCut}
\label{sec:qaoa}

QAOA (\emph{Quantum Approximate Optimization Algorithm}) предложили су
Farhi, Goldstone и Gutmann као хибридни квантно-класични алгоритам за
комбинаторну оптимизацију~\cite{farhi2014}. За наше потребе није
пресудан детаљан профил перформанси QAOA, већ чињеница да он показује
како се \emph{иста} проблемска формулација, добијена у облику
QUBO/Изинговог модела, може користити и у моделу квантних кола.

\textbf{Веза QUBO $\to$ Изинг $\to$ QAOA.}
Сваки QUBO проблем
се супституцијом $x_i = (1-s_i)/2$ директно пресликава у Изингов
Хамилтонијан облика
$H = \sum_i h_i s_i + \sum_{i<j} J_{ij} s_i s_j$,
где спинови $s_i \in \{-1,+1\}$ одговарају кубитима у стањима
$|0\rangle$ и $|1\rangle$.
QAOA затим узима тај Хамилтонијан \emph{директно} као проблемски
оператор $\hat{H}_C$ без икакве додатне трансформације.
Другим речима, QUBO матрица $Q$ потпуно одређује квантно коло QAOA-е:
дијагонални елементи $Q_{ii}$ дају локална поља $h_i$,
а вандијагонални $Q_{ij}$ дају спрезања $J_{ij}$.

\subsection{Од MaxCut-а до проблемског Хамилтонијана}

За MaxCut уводимо спин варијабле $z_v \in \{-1,+1\}$ за сваки чвор
$v \in V$. Ивица $\{u,v\}$ прелази рез ако и само ако $z_u z_v = -1$,
па је величина реза
\[
  C(\mathbf{z}) = \sum_{\{u,v\}\in E} \frac{1-z_u z_v}{2}.
\]
Ова циљна функција директно прелази у проблемски Хамилтонијан
\[
  \hat{H}_C = \frac{1}{2}\sum_{\{u,v\}\in E}\bigl(\hat{I}-\hat{Z}_u\hat{Z}_v\bigr).
\]

Кључна ствар је да је то \emph{иста} енергетска структура коју смо већ
добили из QUBO формулације у поглављу~\ref{sec:qubo}. Зато QAOA није
нова формулација проблема, већ нови начин да се оптимизује већ познати
Хамилтонијан.


\subsection{Зашто је QAOA важан}

У контексту наше дискусије, QAOA је важан из три разлога.
\begin{itemize}[nosep]
  \item Показује да се MaxCut може решавати у моделу квантних кола над
        истим QUBO/Изинговим описом који користе квантни анилери.
  \item Илуструје хибридни образац: квантни део претражује простор
        стања, а класични део управља оптимизацијом.
  \item Пружа природно место за поређење са Z3: QAOA тражи добро
        решење варијационо, док Z3 може егзактно да провери мале
        инстанце и исправност добијеног bitstring-а.
\end{itemize}

Практично ограничење QAOA у NISQ окружењу јесте то што дубина кола,
шум и цена класичне оптимизационе петље брзо постају доминантни.
Зато је у овом раду QAOA пре свега \emph{концептуални мост} између
QUBO формулације и квантног рачунарства заснованог на колима и указујемо на њега 
као пример, како се исти проблем може решавати у различитим моделима рачунања.
